\documentclass[11pt,a4paper]{article}
\usepackage[T1]{fontenc}
\usepackage[utf8]{inputenc}
\usepackage{lmodern}
\usepackage{amsmath,amssymb}
\usepackage[margin=24mm]{geometry}
\usepackage{longtable,booktabs,array,calc}
\usepackage{graphicx}
\usepackage[hidelinks]{hyperref}
\hypersetup{pdftitle={Action, Trajectories and Gaps: A Research Programme},pdfauthor={navstokgap research working notes}}
\newcommand{\tightlist}{\setlength{\itemsep}{0pt}\setlength{\parskip}{0pt}}
\setlength{\emergencystretch}{3em}
\title{Action, Trajectories and Gaps: A Research Programme}
\author{navstokgap research working notes}
\date{}
\begin{document}
\maketitle
\begin{abstract}
We organize the accepted action-scale results around the physical premises
that supply positivity and the counterexamples a selection principle must
exclude. Quantum reconstruction and interacting spectral control are the main
research tracks. The programme identifies their next proof obligations and
keeps dimensional action selection, calibrated observable access and uniform
limits explicit. The maintained source is \texttt{research/PROGRAMME.md};
\texttt{research/STATE.md} supplies the single current priority list.
\end{abstract}
\hypertarget{research-programme-quantum-necessity-and-robust-gap-mechanisms}{%
\section{Research programme: quantum necessity and robust gap
mechanisms}\label{research-programme-quantum-necessity-and-robust-gap-mechanisms}}

Version 72, 2026-09-16. The programme consolidates the accepted results
into \href{../notes/action-scale-obstructions.md}{Classical action
scales: obstructions, conditional bounds and quantum premises}. The
user's clarified 2026-09-14 priority is a mathematical necessity
argument for the Newton/Galileo area--action obstruction; gap work is
supporting. \href{STATE.md}{STATE} owns the current ordered tasks;
\href{STRATEGY.md}{STRATEGY} governs bounded choices and stopping
decisions.

\hypertarget{aim-and-present-conclusion}{%
\subsection{Aim and present
conclusion}\label{aim-and-present-conclusion}}

Derive the need for quantum structure with a positive universal action
parameter from explicit physical premises, and develop a precise
mechanism for a gap that survives the relevant limits. A construction
admitting quantum inputs and a construction selecting those inputs
answer different questions. Newtonian cut refinement remains the
consistency test for proposed action mechanisms; the Millennium
comparison supplies the field-theory target.

The accepted examples provide a substantial set of obstructions to
simple classical selection arguments and several conditional positive
bounds. The physical premise excluding their quiet, contracting or
precisely observed countermodels remains open. Spectral work gives
finite-system criteria, independent-product control and C124's
interacting Ising gap uniform in finite volume at bounded coupling.
Physical-operator identification and further uniform limits remain open.

This version integrates C129/B77: the initial reversible-generator
restriction is proof-audited, with the later interaction-theorem steps
still open. C128/B74 establishes that the exact finite observable repair
of C127 is impossible for the admitted spin preparations. G04 retains
its local Hermitian parent C126 alongside G03's interacting gap C124. No
novelty is asserted. \href{../claims/LEDGER.md}{The ledger} retains
proof and literature status, including the limits of each source audit.

\hypertarget{quantum-track-identify-the-excluded-classical-alternative}{%
\subsection{Quantum track: identify the excluded classical
alternative}\label{quantum-track-identify-the-excluded-classical-alternative}}

Q01's first milestone compares Hardy and Chiribella--D'Ariano--Perinotti
through the exact premise that rules out a classical operational model.
The bounded \href{../notes/quantum-exclusion-premises.md}{premise map}
and B67 audit are complete: Hardy requires continuous reversible
pure-state connectivity; CDP requires pure extensions with essential
uniqueness. Standard finite classical models fail these premises.
Hardy's moving-ball example points to the extra question of an exact
finite operational description closed during transformations. Neither
full reconstruction proof has been audited here. For each selected
passage specify the state space, preparations, transformations,
composition and measurement assumptions. Test the classical alternative
in that same class and mark the step where action units would still need
to enter.

K01 is the supporting fixed-time compatibility test. It asks whether
local measurement descriptions extend to one system-plus-apparatus state
and which assumption about context independence changes the answer. A
primary Kochen--Specker comparison is a source task with explicit
hypotheses. Merely hiding apparatus variables is not enough to exclude
their classical model.
\href{../ideas/I007-static-compatibility.md}{I007} records the
construction.

The \href{../references/batches/B01.md}{B01 companions} provide the
existing Hardy/CDP source routes and limited reading coverage. Full
reconstruction-proof acceptance remains future work. The historical
source-to-model connections are preserved under their own tasks and can
supply a specific premise or example.

\hypertarget{action-selection-obligations}{%
\subsubsection{Action selection
obligations}\label{action-selection-obligations}}

The user's central example is Galileo's falling parabola compared with
the inertial horizontal line. Its enclosed area \(A\) obeys
\((3F/v)A=\tau\Delta E\), with \(\Delta E=F^2\tau^2/(2m)\) the gained
kinetic energy. The matched-endpoint chord lens is half that area;
Kepler's swept sectors are a different observable. The objective is a
physical obstruction to indefinite refinement into distinguishable
classical motion, connected to a positive universal action scale.

\href{../notes/newton-insertion-action.md}{N01} makes a closed
comparison starting with the falling parabola: an explicit return-force
sequence reunites the arms with equal positions and momenta, and its
exact quantum relative phase is \(-2F^2\tau^3/(3m\hbar)\). With C006 it
gives a positive resolution threshold at fixed copies and accuracy. A
hard area cutoff does not follow from the deterministic energy gain
alone; the quantum construction supplies \(\hbar\). This is a supporting
quantum benchmark. The intended innovation is a mathematical necessity
argument from independently justified consistency premises that exclude
zero-action refinement without supplying quantum kinematics. N02 tests
one such proposed principle with an explicit cut state, composition rule
and joint time/position limit. The N01 construction remains exploratory
and supplies no necessity theorem.

\href{../notes/action-unit-dimensional-selection.md}{Q14's dimensional
criteria} consolidate the selection countertests: a floor shared by all
admitted masses and preparations is a dimensionless multiple of an
action-dimensional product of the fixed constants, and an admitted
similarity that rescales the observable forces the floor to zero. The
fixed constants of classical mechanics, electrodynamics and gravitation
admit one mass-independent action,
\(k_e/c=e^2/(4\pi\epsilon_0c)=\alpha\hbar\), and C052's relativistic
Kepler threshold attains it. The open selection premises are a fixed
charge unit and the factor \(1/\alpha\). The criteria remain
exploratory. C130, accepted by user direction after the
\href{../reviews/kepler-collapse-Q14.md}{collapse audit}, records that
C052's plunge threshold and circular energies at angular action
\(\hbar|\kappa|\) are the Dirac Coulomb thresholds and zero-radial-node
levels, and that radial action quantization of C052 gives the Sommerfeld
formula, which is the Dirac spectrum; Planck's constant enters only as
the unit of angular action at the plunge boundary.

\href{../notes/local-detector-coincidences.md}{Q12's local receiver
test} derives an all-gate coincidence-product bound for independent
monotone responses to fixed fractions of one classical pulse, and a
sharp efficiency/balance trade-off. Fluctuating routing, shared
readiness and outcome-based gate selection supply explicit countermodels
when the corresponding premises are removed. This exploratory result
limits a receiver class; its dimensionless statistics supply no action
scale. Shared-preparation models require a separate test involving
independently chosen local settings.

\href{../notes/thermal-receiver-reliability.md}{Q10's thermal receiver
test} combines a dark retention interval with a response deadline. Its
activated-rate model yields an exact barrier-change floor proportional
to temperature and a logarithmic timing/error contrast. A fixed
signal-to-work coupling is needed to turn that into an incident-energy
bound. Kinetic gating exposes that missing premise; temperature and
apparatus parameters remain. This closes arbitrary preloading only
within the stated effective model, without ledger promotion.

\href{../notes/mechanical-interference-action.md}{Q08's mechanical
interference test} constructs an ideal lossless string network with two
paths and matched work-recording receivers. Exact finite-pulse output
energies contain the wave autocorrelation; long coherent pulses give
cosine fringes. Attenuation preserves normalized fringes while energy
and canonical action vanish. The canonical action-to-phase ratio depends
on preparation and impedance; the traveling wave's spacetime Lagrangian
action is zero. This exploratory result tests Q07's physical phase
premise without ledger promotion. Localized event formation and its
energy source remain a distinct receiver question.

\href{../notes/reciprocal-coupling-normalization.md}{Q06's reciprocal
coupling test} constructs conservative exchange with measured
acceleration responses. Reciprocity determines relative inertias,
subject to a network cycle condition; one common energy/action
multiplier survives. This exploratory construction extends Q05's
single-field normalization test to interacting sectors without promoting
a ledger claim. It supplies relative calibration, while absolute
normalization and the quantum phase role remain separate obligations.

C127's \href{../notes/hamiltonian-moment-descent.md}{Hamiltonian test}
now separates microscopic reversibility from operational descent. Two
zero-energy spin preparations share all retained first/cross moments but
an allowed product test distinguishes them after the interaction. The
\href{../notes/hamiltonian-finite-closure.md}{finite-repair proof}
C128/B74 now rules out every finite observable enlargement on a nonzero
time interval for all Borel preparations, including nonlinear updates of
finite expectation vectors. Exact closure requires infinite observable
dimension or a changed physical premise. This settles the named repair
alternative; it selects neither quantum structure nor a universal action
unit. Further hierarchy variants are parked while the missing physical
premises are consolidated.

P06 now consolidates these results into a standalone
\href{../papers/classical-spins-operational-closure.tex}{spin
manuscript},
\href{../out/papers/classical-spins-operational-closure.pdf}{compiled
PDF} and
\href{../out/publications/classical-spins-operational-closure.zip}{PDF-inclusive
bundle}. The \href{../reviews/spin-publication-P06.md}{consistency
review} maps claims to existing audits.
\href{../reviews/spin-publication-B76.md}{B76} supports a worked
conceptual example: classical extensions and transformation-relative
preparation equivalence supply established precedents; the explicit spin
pair and finite-repair proof supply the model-specific presentation. The
manuscript now distinguishes the initial quotient from the expanded
experimental theory. Its action-selection premise remains a separate
physical question.

Q04's \href{../notes/topological-sector-action-selection.md}{topological
field test} constructs a conservative O(3) sector with energy at least
\(4\pi\rho\), while an exact degree-one family of radius \(R\) has
radius-crossing action \(4\pi\rho R/c\to0\). The static energy floor
coexists with massless vacuum waves and translation zero modes. This
exploratory test separates sector exclusion from a physical size or time
scale; Q03's fueled attraction and Q04's geometric restriction have not
supplied a universal action constant.

A candidate must discharge five gates in
\href{ACTION_FIELD_TARGET.md}{the action target}: classical definition
and units; exclusion of zero; convergence in a named topology;
universality across the admitted masses and preparations; and
identification with quantum phase or commutation structure. Define
\(\mathsf h_\varepsilon(t)\) when using the proposed action-field route,
keeping physical time t and resolution epsilon distinct. Observation
duration Delta, volume and \(c^{-1}\) have their own limits.

The \href{../notes/checkerboard-dynamics.md}{checkerboard calculation}
takes a supplied action parameter, amplitudes and measurement rule to a
coherent Dirac limit. It is the Branch A comparison for a Branch B
necessity argument. A state-space reconstruction still needs a
mechanical dynamical identification of its dimensional scale; numerical
calibration is a further empirical task.

\hypertarget{gap-track-interactions-and-uniform-control}{%
\subsection{Gap track: interactions and uniform
control}\label{gap-track-interactions-and-uniform-control}}

Q01's \href{../notes/classical-orientation-closure.md}{orientation
countermodel} (C125/B69) makes the quantum-track boundary explicit:
finite capacity and reversible local closure coexist with classical
separable composition, failed purification and failed subspace
equivalence. The \href{../notes/reversible-interaction-premise.md}{B70
interaction audit} now excludes a continuous reversible nonlocal
extension preserving that minimal composite, using de la Torre et al.'s
Theorem 1. The local ball and Hermitian coordinates already meet the
local-qubit premise. Identical-copy consistency and operational
ancilla/measurement/discard closure underlie the stronger quantum
reconstruction. The full theorem proof remains unaudited here; physical
reversible descent and action normalization remain open. G04's
\href{../notes/ising-hermitian-transfer.md}{Hermitian transfer}
(C126/B71) gives a positive three-site local parent with square-root
Gibbs ground state and energy-unit gap \(Ka[1-\tanh(2b)]\). Both action
constant K and the physical clock identification are supplied. This
resolves the finite-volume operator construction while retaining the
independent physical-dynamics premise.

\href{../notes/finite-depth-spin-gap.md}{G05's direct Hamiltonian} works
through an established cluster chain obtained by a depth-two
controlled-Z circuit. It has three-site interactions, an entangled
ground state and exact energy gap J uniform in chain length at fixed
lattice spacing and coupling. Endpoint terms control ground-state
uniqueness. This gives a direct quantum energy operator without a
sampling-clock identification, with quantum kinematics and the action
factor for physical time supplied. Robustness under a local noncommuting
perturbation is a separate test; no new ledger claim is promoted.

\href{../notes/low-dimensional-mass-gap.md}{G07's dictionary note}
answers the user's 2026-09-16 question of how h \textgreater{} 0 is
analogous to the mass gap. The floor/unit theorem of Q14 applies
verbatim to both observables: a positive finite floor is a dimensionless
multiple of a unit formed from the fixed constants, and classical
Yang--Mills in four dimensions and classical mechanics with a mass and a
force both lack the unit. The note proves, with \(\hbar\), \(g\) and
\(m\) explicit, that Yang--Mills quantum mechanics (the zero-momentum
sector of the torus theory) has a purely discrete spectrum and gap
\(\delta_1\hbar^{4/3}g^{2/3}m^{-2/3}\) which vanishes when either the
canonical or the Lie-algebra commutator is removed; that free and
abelian fields inherit their gapless classical dispersion; that
two-dimensional Yang--Mills on a circle has a gap
\(\tfrac12g^2\hbar^2cL\,C_2\) diverging with the volume; and that
scale-invariant quantum mechanics obtains its scale from the
self-adjointness requirement. The finite-volume zero-mode gap
\(\delta_1g^{2/3}\hbar c/L\) exhibits Jaffe--Witten's uniformity problem
in a solved sector. Consequence: C124/C126 and G05 are supplied-scale
gaps; a generated gap needs two scale-free non-commuting structures; the
analogy covers the structure-to-scale stage of the necessity question
and points N02 to a consistency-forced scale. Exploratory; sources in
B78; no promotion.

G01/G02 establish the finite reversible-generator estimate

\[\gamma\ge\frac{\alpha}{S},\qquad
S=\sum_a\langle f_a,(-Q)^{-1}f_a\rangle,\]

where alpha is the lower observable-frame bound on centered states. One
positive measured plateau can miss an arbitrarily slow mode; exact state
identification can also become poorly conditioned. Independent product
dynamics control mixed modes and give the minimum constituent gap when
each constituent clock is retained (C041--C046).

G03 now supplies an explicit interacting estimate: the periodic
zero-field Ising heat-bath chain has full gap \(a[1-\tanh(2b)]\) for
every \(N\ge3\). Here \(a\) is the per-site refresh rate and \(b\ge0\)
the dimensionless nearest-neighbour coupling. Hamming contraction
controls all modes and magnetization attains the bound. Bounded coupling
and a clock floor give a uniform finite-volume gap; finite range alone
permits joint coupling/size closure, and a fixed total clock adds a
\(1/N\) slowdown. The \href{../notes/interacting-ising-gap.md}{C124
proof} replaces tensorization without an extensive full-frame
susceptibility sum. B68 records a bounded model/method audit, including
unavailable Glauber formula images. Further interaction variants are
parked until a named operator or limit dependency requires them. See
also \href{../notes/susceptibility-gap.md}{the susceptibility proof}.

A Yang--Mills transfer additionally requires its physical quantum
Hamiltonian, gauge-invariant observables, continuum and infinite-volume
construction and quantum axioms. An auxiliary sampling generator's
relaxation rate changes under clock rescaling while its invariant law
stays fixed. The \href{../notes/comparison-and-bridges.md}{comparison
note} and
\href{../notes/millennium-problem-definitions.md}{official-target
digest} preserve these distinctions. A finite-box decay estimate or a
toy Hessian gap supplies only its own operator statement.

\hypertarget{consolidated-evidence-and-retained-supporting-models}{%
\subsection{Consolidated evidence and retained supporting
models}\label{consolidated-evidence-and-retained-supporting-models}}

\begin{longtable}[]{@{}
  >{\raggedright\arraybackslash}p{(\columnwidth - 4\tabcolsep) * \real{0.3333}}
  >{\raggedright\arraybackslash}p{(\columnwidth - 4\tabcolsep) * \real{0.3333}}
  >{\raggedright\arraybackslash}p{(\columnwidth - 4\tabcolsep) * \real{0.3333}}@{}}
\toprule\noalign{}
\begin{minipage}[b]{\linewidth}\raggedright
Mechanism tested
\end{minipage} & \begin{minipage}[b]{\linewidth}\raggedright
Accepted consequence
\end{minipage} & \begin{minipage}[b]{\linewidth}\raggedright
Premise still requiring selection
\end{minipage} \\
\midrule\noalign{}
\endhead
\bottomrule\noalign{}
\endlastfoot
Continuous variation and scale contraction & C002/C056--C058 give zero
action infima in specified classes & Restriction excluding the allowed
variations or low excitation \\
Positive orbit cost & C060--C061 give a sharp force/speed-floor bound &
Physical origin and universality of the excitation floor \\
Fluctuation and independent composition & C035--C036 give a common
coefficient using a positive reference & Preparation equivalence and
origin of reference fluctuations \\
Prepared apparatus and records & C068--C123 give recovery and ambiguity
under explicit data restrictions & Irreducible, universal
information/preparation constraints \\
Coherent paths & C039--C040 give continuum and measured-cut comparisons
& Necessity of amplitudes, measurement rule and action parameter \\
Observable response and interaction & C041--C046 give gap criteria; C124
gives an interacting uniform finite-volume gap & Origin of
coupling/clock control, physical operator identification and further
limits \\
\end{longtable}

The \href{../out/papers/action-scale-obstructions.pdf}{consolidated
paper} supplies short arguments and links to the full proofs. The
receiver/calibration compilation retains technical details, including
R33's open curvature test. R33 is parked: its next coefficient currently
changes no selected quantum or gap premise. It can return for a named
dependency, proof correction or an explicit user request.

A19's peak-excitation estimate, A14's renewal-preparation model, R02's
repaired finite-propagation drafts and M03's spectral laboratory remain
supporting work. They become selected work when their possible outcomes
change a main premise decision. New solvable subcases do not
automatically displace the main tracks.

\hypertarget{cut-refinement-and-historical-context}{%
\subsection{Cut refinement and historical
context}\label{cut-refinement-and-historical-context}}

The \href{CUT_POINT_TARGET.md}{cut-point target} collects the existing
consistency tests. Constant-force chord error closes under refinement;
fixed Gaussian bridge laws preserve their supplied fluctuation
parameter; velocity and receiver memory distinguish retained-state cuts
from interventions. Joint path/action limits require their stated
topology and preparation changes. These results constrain proposed
mechanisms without selecting quantum weights.

Newton's geometric arguments, the cone/arrow proposals, the Classical
Scholia and the H10 static-composition readings remain source-backed
research inputs. \href{../references/IDEA_BRIDGES.md}{The bibliography
bridge index} routes a named task to a construction, premise or exact
passage. H11 strengthens source witnesses when selected; it does not
resume merely because an old source card calls it next. The supplied
\href{../references/nota_unificada.md}{bibliography note} remains
intact, with its
\href{../references/nota_unificada-seleccion-matematica.md}{separate
mathematical selection}.

\hypertarget{publication-development}{%
\subsection{Publication development}\label{publication-development}}

The user has selected publication preparation. P05 extracts C050--C051
into \href{../out/papers/same-collisions-different-transport.pdf}{Same
collisions, different transport}, a standalone teaching draft with a
\href{../out/publications/same-collisions-different-transport.zip}{PDF-inclusive
source bundle}. B75 finds a close Brownian random-versus-regular
preparation precedent; the exact Hamiltonian matched-rate comparison is
unmatched only within the bounded search. No new theorem or novelty
claim is introduced. The spin draft consolidates C125--C128, including
the completed finite-repair obstruction. P07 completes accessibility and
contribution edits for both drafts; the
\href{../reviews/publication-readiness-P07.md}{author decision sheet}
recommends the gas paper first and records the remaining author and
destination choices. STATE owns ordering; journal submission remains a
separate authorized action.

\hypertarget{acceptance-cost-and-consolidation}{%
\subsection{Acceptance, cost and
consolidation}\label{acceptance-cost-and-consolidation}}

Use written proofs and bounded source review. No numerical or symbolic
verification scripts may be created or run in Python or another
language. Historical scripts and outputs remain records.
Document/source/integrity Python is allowed. Delegation is selective and
sequential under \href{../agents/PROTOCOL.md}{the protocol}, with one
Sol/Luna worker and explicit supported effort, never ultra.

Before each step, record its target obligation, the decision changed by
its possible outcomes, its deliverable/stop and why existing results are
insufficient. After two consecutive steps in a proof family, compare its
continuation with the other main track. The coordinator makes that
choice without routine user approval. Correct substantive errors when
found and preserve explicit user scope. A completed local calculation
need not spawn another calculation.

A new result receives written proof review and a bounded librarian
audit. Consolidation reuses those audits with an evidence map; it does
not imply a fresh prior-art search or automatic novelty. Keep STATE
short and reconcile its queue with TASKS and the latest handoff. Dated
milestones live in \href{STATE-history-2026-09-12.md}{the state history}
and individual handoffs.

Run \texttt{make\ check}, \texttt{make\ papers} and
\texttt{git\ diff\ -\/-check}; inspect changed PDF pages and preserve
unrelated worktree changes. Relevant status/results are committed and
pushed under standing authorization. External submissions,
correspondence and paid services need separate direction. Publication
claims must reflect the accepted proofs and the actual literature
coverage.
\end{document}
